\documentclass[../../main.tex]{subfiles}

\begin{document}

\chapter[Finite Automata]{Finite Automata: NFA, DFA, Konversi,\\ dan Minimalisasi}
\label{chap:finite-automata}

\begin{outcomebox}
\textbf{Sub-CPMK:} \subcpmk{2.2}, \subcpmk{2.3}\\[3pt]
Setelah menyelesaikan bab ini, mahasiswa mampu:
\begin{enumerate}[label=\arabic*.,leftmargin=*]
  \item mendefinisikan DFA dan NFA beserta cara kerjanya;
  \item menghitung epsilon closure dan menyimulasikan NFA;
  \item mengonversi NFA menjadi DFA dengan subset construction;
  \item meminimalkan DFA dengan penyempurnaan partisi;
  \item menghubungkan otomata hingga dengan pengenalan token.
\end{enumerate}
\end{outcomebox}

\section*{Materi Pokok}
\begin{itemize}[leftmargin=*]
  \item DFA: definisi, tabel transisi, dan simulasi.
  \item NFA: nondeterminisme dan epsilon closure.
  \item Konversi NFA ke DFA dengan subset construction.
  \item Minimalisasi DFA dan partisi ekuivalensi.
  \item Dari ekspresi reguler ke otomata dan token recognition.
\end{itemize}

\subfile{section-01}
\subfile{section-02}
\subfile{section-03}
\subfile{section-04}
\subfile{section-05}

\section*{Contoh}
\addcontentsline{toc}{section}{Contoh}

\begin{contoh}[Menelusuri DFA]
Gunakan DFA pada \figref{fig:ch06-dfa-genap} untuk memeriksa string
\(abba\). State awal \(q_0\). Membaca \(a\): pindah ke \(q_1\).
Membaca \(b\): tetap di \(q_1\). Membaca \(b\): tetap di \(q_1\).
Membaca \(a\): pindah ke \(q_0\). State akhir \(q_0\) adalah state
penerima, sehingga \(abba\) diterima. Jumlah \(a\) pada string adalah
dua (genap), sesuai definisi bahasa.
\end{contoh}

\begin{contoh}[Epsilon closure bertingkat]
Misalkan NFA memiliki transisi
\(q_0 \xrightarrow{\varepsilon} q_1\),
\(q_1 \xrightarrow{\varepsilon} q_2\),
\(q_2 \xrightarrow{a} q_3\),
\(q_3 \xrightarrow{\varepsilon} q_4\).
Maka
\[
\varepsilon\text{-closure}(\{q_0\}) = \{q_0,q_1,q_2\}.
\]
Setelah membaca \(a\), himpunan state menjadi
\(\{q_3\}\), lalu epsilon closure-nya menjadi
\(\{q_3,q_4\}\). Jika \(q_4\) adalah state penerima, string \(a\)
diterima.
\end{contoh}

\begin{contoh}[Subset construction]
Untuk NFA dengan transisi
\(q_0 \xrightarrow{a} \{q_0,q_1\}\),
\(q_0 \xrightarrow{b} \{q_0\}\),
\(q_1 \xrightarrow{b} \{q_2\}\), dan \(q_2\) penerima:
\begin{center}
\small
\begin{tabular}{@{}c l c c c@{}}
\toprule
\textbf{State DFA} & \textbf{Himpunan NFA} &
\textbf{\(a\)} & \textbf{\(b\)} & \textbf{Penerima} \\
\midrule
\(S_0\) & \(\{q_0\}\) & \(S_1\) & \(S_0\) & tidak \\
\(S_1\) & \(\{q_0,q_1\}\) & \(S_1\) & \(S_2\) & tidak \\
\(S_2\) & \(\{q_0,q_2\}\) & \(S_1\) & \(S_0\) & ya \\
\bottomrule
\end{tabular}
\end{center}
State penerima adalah \(S_2\), karena memuat \(q_2\). DFA hasilnya
memiliki tiga state.
\end{contoh}

\begin{contoh}[Minimalisasi]
Diberikan DFA dengan state \(\{0,1,2,3\}\), state penerima
\(\{3\}\), dan transisi:
\[
0 \xrightarrow{a} 1,\quad
0 \xrightarrow{b} 2,\quad
1 \xrightarrow{a} 1,\quad
1 \xrightarrow{b} 3,\quad
2 \xrightarrow{a} 1,\quad
2 \xrightarrow{b} 2,\quad
3 \xrightarrow{a} 1,\quad
3 \xrightarrow{b} 2.
\]
Partisi awal \(\{0,1,2\}\) dan \(\{3\}\). Tabel berikut mencatat
setiap iterasi.

\begin{center}
\small
\begin{tabular}{@{}c l l@{}}
\toprule
\textbf{Iterasi} & \textbf{Partisi} & \textbf{Alasan} \\
\midrule
0 & \(\{0,1,2\}, \{3\}\) & penerima versus bukan penerima \\
1 & \(\{0,1\}, \{2\}, \{3\}\) &
\(1 \xrightarrow{b} 3\), sedangkan \(2 \xrightarrow{b} 2\) \\
2 & \(\{0\}, \{1\}, \{2\}, \{3\}\) &
\(0 \xrightarrow{b} 2\), sedangkan \(1 \xrightarrow{b} 3\) \\
\bottomrule
\end{tabular}
\end{center}

Hasil akhirnya empat grup; DFA sudah minimal.
\end{contoh}

\begin{contoh}[Token recognition]
Misalkan DFA gabungan memiliki state penerima untuk identifier dan
angka. Pada masukan \code{abc123}, DFA membaca seluruh string dan
berhenti di state penerima identifier. Pada masukan \code{123abc},
DFA membaca \code{123} sebagai angka, lalu berhenti karena transisi
berikutnya tidak ada. Lexer mengeluarkan token angka \code{123} dan
memproses \code{abc} pada pemanggilan berikutnya. Maximal munch
memastikan \code{123} tidak dipecah menjadi tiga token digit.
\end{contoh}

\section*{Aktivitas Pembelajaran dan Praktikum}
\addcontentsline{toc}{section}{Aktivitas Pembelajaran dan Praktikum}

\begin{table}[H]
\centering
\caption{Rencana aktivitas Bab VI}
\label{tab:ch06-aktivitas}
\begin{tabularx}{\textwidth}{@{}llX@{}}
\toprule
\textbf{Sesi} & \textbf{Bentuk} & \textbf{Kegiatan} \\
\midrule
1 & Ceramah & DFA, NFA, dan epsilon closure \\
2 & Latihan & Menelusuri string pada DFA dan NFA \\
3 & Lokakarya & Subset construction, tabel konversi, dan minimalisasi \\
4 & Praktikum & P1-B: simulator NFA/DFA C dan uji ekivalensi \\
\bottomrule
\end{tabularx}
\end{table}

\subsection*{P1-B: Simulator NFA/DFA dan Tabel Konversi}
\textbf{Tujuan.} Membuat tabel konversi NFA ke DFA, tabel minimalisasi,
dan simulator NFA/DFA dalam C, lalu menguji bahwa keduanya menerima
bahasa token yang sama.

\textbf{Langkah kerja.}
\begin{enumerate}[label=\arabic*.,leftmargin=*]
  \item Untuk setiap kategori token pada \file{spec/tokens.md}, gambar
        NFA yang mengenali kategori tersebut.
  \item Lakukan subset construction dan isi tabel konversi: state DFA,
        himpunan state NFA, transisi setiap simbol, dan status
        penerima.
  \item Minimalisasi DFA dan isi tabel partisi: iterasi, partisi, serta
        alasan setiap pemisahan.
  \item Implementasikan simulator NFA dan DFA dalam C pada
        \file{nfa\_sim.c} dan \file{dfa\_sim.c}.
  \item Uji kedua simulator dengan 30 kasus uji dari P1-A.
  \item Pastikan NFA, DFA hasil konversi, dan DFA minimal menerima
        bahasa token yang sama.
  \item Bandingkan hasilnya dengan \file{fa\_tool.py}; catat perbedaan
        perilaku sebagai pembanding, bukan sebagai artefak utama.
\end{enumerate}

\textbf{Definisi selesai.}
\begin{checklist}
  \item Tabel konversi NFA ke DFA lengkap dan konsisten.
  \item Tabel minimalisasi memuat alasan setiap pemisahan.
  \item Simulator NFA dan DFA dalam C berjalan.
  \item 30 kasus uji menghasilkan keputusan yang sama pada ketiga
        automata.
  \item Perbedaan dengan \file{fa\_tool.py} terdokumentasi.
\end{checklist}

P1-B menjadi bagian Praktikum 1 (12\%) yang dilaporkan pada Bab VII.

\section*{Latihan dan Refleksi}
\addcontentsline{toc}{section}{Latihan dan Refleksi}

\subsection*{Latihan}
\begin{enumerate}[label=\arabic*.,leftmargin=*]
  \item Gambarkan DFA untuk bahasa string atas \(\{a,b\}\) yang
        memuat jumlah \(b\) kelipatan tiga.
  \item Hitung epsilon closure dari \(\{q_0\}\) pada NFA dengan
        transisi \(q_0 \xrightarrow{\varepsilon} q_1\),
        \(q_1 \xrightarrow{\varepsilon} q_2\), dan
        \(q_2 \xrightarrow{a} q_3\).
  \item Lakukan subset construction untuk NFA pada Latihan 2 dan isi
        tabel konversi: state DFA, himpunan state NFA, transisi, dan
        status penerima.
  \item Minimalisasi DFA yang dihasilkan dan isi tabel partisi:
        iterasi, partisi, serta alasan setiap pemisahan.
  \item Uji ekivalensi NFA, DFA hasil konversi, dan DFA minimal dengan
        lima string; jelaskan hasilnya.
  \item Jelaskan mengapa jumlah state DFA hasil subset construction
        paling banyak \(2^{|Q|}\).
  \item Rancang NFA untuk \((ab \mid ba)^{*}\).
  \item Terapkan konstruksi Thompson pada \((a \mid b)^{*}abb\).
  \item Mengapa lexer memerlukan rollback pada maximal munch?
\end{enumerate}

\subsection*{Refleksi}
\begin{reflectionbox}
\begin{enumerate}[label=\arabic*.,leftmargin=*]
  \item Apakah Anda dapat menjelaskan perbedaan DFA dan NFA tanpa
        melihat catatan?
  \item Bagian mana dari subset construction yang paling sulit
        dikerjakan manual?
  \item Mengapa minimalisasi penting untuk ukuran tabel lexer?
\end{enumerate}
\end{reflectionbox}

\section*{Asesmen}
\addcontentsline{toc}{section}{Asesmen}

Asesmen Bab VI mencakup Tugas Teori (8\%) dan bagian P1-B. Fokusnya
adalah kemampuan mengonversi dan meminimalkan otomata.

\begin{table}[H]
\centering
\caption{Instrumen asesmen Bab VI}
\label{tab:ch06-asesmen}
\begin{tabularx}{\textwidth}{@{}lXl@{}}
\toprule
\textbf{Instrumen} & \textbf{Indikator} & \textbf{Bobot} \\
\midrule
Tugas Teori & DFA/NFA, subset construction, minimalisasi & 8\% \\
P1-B & DFA token dan prototipe scanner & bagian P1 \\
Kuis & Epsilon closure dan simulasi & formatif \\
\bottomrule
\end{tabularx}
\end{table}

Kriteria kelulusan: mahasiswa dapat mengonversi NFA menjadi DFA
secara manual, meminimalkan DFA, dan menjelaskan hubungannya dengan
pengenalan token.

\section*{Checklist Pencapaian Kompetensi}
\addcontentsline{toc}{section}{Checklist Pencapaian Kompetensi}

\begin{checklist}
  \item Saya dapat mendefinisikan DFA dan NFA.
  \item Saya dapat membuat tabel transisi DFA.
  \item Saya dapat menghitung epsilon closure.
  \item Saya dapat menyimulasikan NFA dengan himpunan state.
  \item Saya dapat mengimplementasikan simulator NFA dan DFA dalam C.
  \item Saya dapat melakukan subset construction dan mengisi tabel
        konversi.
  \item Saya dapat meminimalkan DFA dan mengisi tabel partisi.
  \item Saya dapat menguji bahwa NFA, DFA hasil konversi, dan DFA
        minimal menerima bahasa token yang sama.
  \item Saya dapat menggabungkan beberapa token menjadi satu DFA.
  \item Saya sudah menyelesaikan P1-B.
\end{checklist}

\section*{Rangkuman}
\addcontentsline{toc}{section}{Rangkuman}

Pada bab ini telah dibahas:
\begin{itemize}[leftmargin=*]
  \item DFA, tabel transisi, fungsi transisi diperluas, dan simulasi;
  \item NFA, nondeterminisme, epsilon closure, dan simulasi;
  \item subset construction untuk mengonversi NFA ke DFA;
  \item minimalisasi DFA dengan penyempurnaan partisi;
  \item konstruksi Thompson dan penggabungan token untuk lexer.
\end{itemize}

Bab VII akan mengimplementasikan seluruh teori ini menjadi lexer yang
bekerja. Kita akan menulis lexer tangan dalam C, lexer Flex, dan
driver yang menghasilkan aliran token untuk parser.

\section*{Evaluasi}
\addcontentsline{toc}{section}{Evaluasi}

\begin{enumerate}[label=\arabic*.,leftmargin=*]
  \item Definisikan DFA dan NFA, lalu sebutkan perbedaannya.
  \item Gambarkan DFA untuk \((a \mid b)^{*}abb\).
  \item Hitung epsilon closure dari \(\{q_0\}\) pada NFA yang
        diberikan.
  \item Lakukan subset construction pada NFA tersebut.
  \item Minimalisasi DFA hasil konversi.
  \item Jelaskan konstruksi Thompson untuk \(r^{*}\).
  \item Mengapa lexer perlu menyimpan posisi terakhir yang cocok?
  \item Jelaskan hubungan antara DFA minimal dan ukuran tabel lexer.
\end{enumerate}

\subsection*{Kunci dan Rubrik Singkat}
\begin{itemize}[leftmargin=*]
  \item Soal 1--3: ketepatan definisi dan perhitungan (70\%),
        penjelasan (30\%).
  \item Soal 4--5: kebenaran konversi (60\%), minimalisasi (40\%).
  \item Soal 6--8: ketepatan konsep (60\%), contoh (40\%).
\end{itemize}


\ifSubfilesClassLoaded{%
  \bibliographystyle{apalike}%
  \bibliography{../../references}%
}{}

\end{document}
